\documentclass[11pt]{article}

% Packages

% Paragraph formatting
\usepackage{setspace}
\setstretch{1.15}

% Remove paragraph indentation
\setlength{\parindent}{0pt}

% Add vertical spacing between paragraphs
\setlength{\parskip}{6pt}

\usepackage[margin=1in]{geometry}
\usepackage{graphicx}
\usepackage{booktabs}
\usepackage{float}
\usepackage{amsmath}
\usepackage{hyperref}
\usepackage[authoryear,round]{natbib}
% Document information
\title{Momentum Has Its Moments: Replication and Extension Using
EWMA and GARCH-Based Volatility Management}
\author{David Poussard (dp3372), Yihang Guo (yg3083), Rong Ji (rj2754)}
\date{December 2026}

\begin{document}

\maketitle


\section{Introduction}
\label{sec:introduction}


Momentum is one of the most extensively documented
anomalies in financial markets. In their seminal study,
\citet{jegadeesh1993returns} demonstrate that strategies
buying stocks with strong past performance and selling
stocks with weak past performance generate significant
positive returns over intermediate investment horizons.
 Stocks that have performed relatively well over
the preceding months tend to continue outperforming stocks with
poor recent performance, giving rise to the widely studied
winner-minus-loser (WML) investment strategy. Despite its
historically attractive average returns, momentum is also
characterized by substantial downside risk. In particular,
momentum strategies can experience severe losses during abrupt
market reversals, when previously underperforming stocks recover
sharply and recent winners underperform.


\citet{barroso2015momentum}, in
\textit{Momentum Has Its Moments}, demonstrate that
momentum risk varies considerably over time and that a relatively simple volatility-management
strategy can substantially improve its risk-adjusted
performance. Their approach adjusts momentum exposure inversely
with an estimate of expected volatility, reducing positions
during turbulent periods and increasing exposure when predicted
volatility is relatively low. Using a six-month rolling window
of daily momentum returns to estimate volatility, they document
a substantial improvement in the strategy's Sharpe ratio,
alongside reductions in extreme losses and other undesirable
features of the return distribution.

These findings raise a natural question: if volatility
management improves momentum performance, can a more
sophisticated volatility forecasting model generate further
benefits? The historical estimator employed by Barroso and
Santa-Clara relies on a fixed window of past squared returns.
Although straightforward and transparent, this approach does
not explicitly model the persistence and clustering of
conditional volatility. Generalized autoregressive conditional
heteroskedasticity (GARCH) models offer an alternative by
allowing expected volatility to evolve dynamically in response
to recent return shocks and past conditional variance.



This paper revisits the volatility-managed momentum
strategy proposed by Barroso and Santa-Clara (2015)
through a replication and two extensions. First, we
reconstruct the momentum portfolio and evaluate
whether scaling exposure using historical volatility
reproduces the principal economic findings of the
original study. Second, we extend the evaluation
through December 2025 to examine whether the
benefits of volatility management persist in a
more recent market environment. Third, we
investigate whether alternative conditional
volatility estimators improve forecasting
accuracy and portfolio performance.

Our methodological extension compares three approaches to forecasting momentum volatility: a 126-day historical estimator, an exponentially weighted moving average (EWMA) estimator, and a rolling GARCH(1,1) model. The historical estimator serves as the benchmark, while EWMA and GARCH allow volatility estimates to respond more dynamically to recent market conditions. We evaluate the resulting strategies using annualized returns, realized volatility, Sharpe ratios, maximum drawdowns, and average portfolio exposure. We also compare out-of-sample volatility forecast accuracy and test differences in squared forecast loss using heteroskedasticity- and autocorrelation-consistent standard errors.

Beyond forecasting accuracy and gross portfolio performance, we investigate whether the economic benefits of volatility management survive realistic implementation constraints. Specifically, we incorporate stock-level proportional transaction costs through a self-financing portfolio simulation and examine sensitivity to leverage caps and financing spreads. Finally, we employ a paired circular moving-block bootstrap to assess the statistical significance of differences in net Sharpe ratios across strategies, accounting for serial dependence and examining robustness to alternative block lengths and leverage constraints.

Our results support the central economic
finding of Barroso and Santa-Clara (2015):
volatility management substantially improves
the risk-adjusted performance of momentum.
During the historical evaluation period,
the unmanaged strategy achieves a Sharpe
ratio of 0.510, compared with 0.963 for
historical volatility management, 1.017
for EWMA, and 0.946 for GARCH. Although
unmanaged momentum performance deteriorates
considerably during the 2012--2025
extension, all three managed strategies
continue to deliver higher realized
Sharpe ratios.

The forecasting results provide a more
nuanced comparison of the conditional
volatility models. Both EWMA and GARCH
reduce forecast errors relative to the
historical estimator, particularly in
the historical sample. EWMA achieves
the lowest historical forecast RMSE,
while GARCH has a marginal advantage
in the extension period. However,
pairwise tests do not establish a
statistically significant difference
in squared forecast loss between
EWMA and GARCH in either period.

Robustness analyses further indicate
that the results do not depend
exclusively on a single conditional
volatility specification. EWMA
performance remains relatively stable
across moderate changes in its decay
factor, while the GARCH analysis
examines alternative innovation
distributions and rolling estimation
windows.

Taken together, these findings
suggest that the principal economic
benefit arises from adjusting
momentum exposure to changing
volatility conditions, rather than
from selecting a uniquely superior
forecasting model. The relative
simplicity and competitive performance
of EWMA also illustrate that greater
model complexity does not necessarily
translate into better realized
investment outcomes.




\section{Data and Portfolio Construction}
\label{sec:data}

\subsection{Data Sources}

This study uses two complementary sources of U.S. equity
returns to replicate and extend the findings of Barroso
and Santa-Clara (2015).

For the replication of the original study, we obtain
value-weighted momentum decile portfolio returns from
Kenneth French's Data Library. We construct the
winner-minus-loser (WML) portfolio as the return
difference between the highest- and lowest-momentum
deciles. This replication covers July 1927 through
December 2011 and provides a benchmark for validating
our implementation of the original volatility-management
strategy.

For our principal extension, we independently reconstruct
a WML momentum portfolio using stock-level return data
from the Center for Research in Security Prices (CRSP),
accessed through Wharton Research Data Services (WRDS).
This reconstruction allows us to apply the historical
volatility estimator and the GARCH(1,1) model to an
identical underlying return series. We also compare the
reconstructed portfolio with the published momentum
factor from Kenneth French's Data Library as an external
validation benchmark.

The extension analysis is divided into two periods.
The historical comparison covers August 1929 through
December 2011, while the subsequent evaluation covers
January 2012 through December 2025. All three strategies
within each period are evaluated on the same monthly
observations, ensuring that differences in performance
are not driven by unequal sample lengths.



\subsection{Momentum Portfolio}
\label{subsec:momentum}

We reconstruct a cross-sectional winner-minus-loser (WML)
momentum strategy using individual U.S. equity returns from CRSP.
Our portfolio formation procedure follows the general framework
of the momentum factor methodology documented in Kenneth
French's Data Library.

At the end of each month $t-1$, stocks are ranked according
to their cumulative returns from months $t-12$ through $t-2$.
The most recent month is excluded to avoid short-term reversal
effects. For stock $i$, the momentum signal used to form the
portfolio held during month $t$ is

\begin{equation}
MOM_{i,t}
=
\prod_{k=2}^{12}(1+r_{i,t-k})-1,
\end{equation}

where $r_{i,t-k}$ denotes the monthly return of stock $i$
at lag $k$. The ranking period therefore contains eleven
calendar months, with a one-month gap before the holding period.

The eligible universe consists primarily of ordinary common
shares (CRSP historical share codes 10 and 11) listed on
the NYSE, AMEX, or NASDAQ. Stocks must have positive
market capitalization at formation, a valid momentum signal,
and an eligible historical price observation. Securities
already delisted at the formation date are excluded.

At each formation date, we calculate momentum breakpoints
using eligible NYSE stocks. The 10th through 90th percentiles
define ten momentum portfolios, and eligible stocks from
all three exchanges are assigned to the corresponding deciles.

Within each decile, stocks are weighted by their market
capitalization at the end of the formation month:

\begin{equation}
w_{i,t-1}^{(d)}
=
\frac{MCAP_{i,t-1}}
{\sum_{j \in \mathcal{P}_{d,t}} MCAP_{j,t-1}},
\end{equation}

where $\mathcal{P}_{d,t}$ is the set of stocks assigned
to decile $d$ for holding month $t$.

The momentum strategy takes a long position in the
highest-momentum decile and a short position in the
lowest-momentum decile:

\begin{equation}
r_t^{WML}
=
r_t^{P10}-r_t^{P1}.
\end{equation}

Portfolios are reconstituted monthly. Within each holding
month, daily portfolio returns are calculated by tracking
the evolution of the value-weighted stock positions,
allowing portfolio weights to drift with realized returns.
The resulting daily reconstructed WML series is used to
estimate and forecast volatility. Monthly WML returns are
constructed separately from the corresponding monthly stock
portfolio returns; they are not obtained by compounding the
daily long-minus-short factor returns.

Our reconstruction explicitly processes recorded delisting
returns and prevents their double-counting. Because some
historical CRSP observations contain incomplete return or
terminal-payoff information, we also implement an explicitly
identified missing-data treatment. The primary reconstructed
daily series retains available numeric return components
and assigns zero increments to otherwise unresolved
components. These assumptions are part of the reconstruction
procedure and should not be interpreted as observed returns
or evidence that the affected securities were tradable.

For the most recent observations, where legacy security
classification fields are unavailable, we apply a
separately identified eligibility screen based on CRSP's
CIZ security attributes. This screen approximates the
historical common-stock universe but is not assumed
to be exactly equivalent to the legacy classification.


\subsection{Sample Construction}

Our historical and extension samples use the same underlying
portfolio construction and volatility-management framework.
This ensures that differences in empirical performance are not
driven by changes in the definition of the momentum strategy.

The historical analysis begins in August 1929 because the
rolling GARCH specification requires an initial estimation
window of 756 daily observations. Restricting the comparison
to dates for which all strategies have available volatility
forecasts ensures a consistent evaluation sample.

The extension period begins in January 2012, immediately
following the historical sample, and continues through
December 2025. Both periods are evaluated using monthly
portfolio returns, while the competing volatility forecasts
are generated from daily data.


\section{Volatility Management Methodology}
\label{sec:methodology}

\subsection{Volatility-Managed Momentum}

Following the framework of Barroso and Santa-Clara (2015),
we investigate whether dynamically adjusting exposure to
momentum based on predicted volatility improves its
risk-adjusted performance.

Let $r_{t+1}^{WML}$ denote the return on the unscaled
winner-minus-loser momentum portfolio in month $t+1$.
The return on the volatility-managed strategy is

\begin{equation}
r_{t+1}^{VM}
=
L_t r_{t+1}^{WML},
\label{eq:managed_return}
\end{equation}

where $L_t$ is the exposure selected at the end of month
$t$, using only information available at that time.

Exposure is determined by

\begin{equation}
L_t
=
\frac{\sigma_{\mathrm{target}}}
{\widehat{\sigma}_{t+1|t}},
\label{eq:vol_target}
\end{equation}

where $\widehat{\sigma}_{t+1|t}$ is the forecast of
momentum volatility for the following month, and
$\sigma_{\mathrm{target}}$ is the desired monthly
volatility level.

We set the annualized volatility target to 12\%,
corresponding to a monthly target of

\begin{equation}
\sigma_{\mathrm{target}}
=
\frac{0.12}{\sqrt{12}}.
\end{equation}

This rule reduces exposure when predicted volatility
is high and increases exposure when predicted volatility
is low. The exposure multiplier is not capped, so
the strategy can employ leverage when forecast
volatility falls below the target.


We compare three approaches to forecasting volatility:
a rolling historical estimator, an exponentially weighted
moving average (EWMA) estimator, and a GARCH(1,1)
conditional volatility model.


\subsection{Historical Volatility Estimator}

Our benchmark follows the historical volatility
estimation approach used in the original study.
At the end of each month $t$, we calculate
volatility using the most recent 126 daily
momentum returns.

Let $r_d^{WML}$ denote a daily momentum return,
and let $\mathcal{D}_t$ contain the 126 trading
days ending at month-end $t$. The forecast is

\begin{equation}
\widehat{\sigma}_{t+1|t}^{HIST}
=
\sqrt{
21 \times
\frac{1}{126}
\sum_{d \in \mathcal{D}_t}
(r_d^{WML})^2
}.
\label{eq:hist_vol}
\end{equation}

The factor of 21 converts the daily second-moment
estimate to an approximate monthly volatility
measure, assuming 21 trading days per month.

The estimator uses squared returns without
subtracting the sample mean, consistent with
the volatility-management specification
implemented in our replication.

\subsection{GARCH(1,1) Volatility Forecast}


Our principal extension replaces the historical
volatility estimate with a conditional volatility
forecast obtained from a GARCH(1,1) model.
Building on the ARCH framework introduced by
\citet{engle1982arch}, the GARCH model developed
by \citet{bollerslev1986garch} captures volatility
clustering through the dependence of conditional
variance on past squared innovations and past
conditional variance.


Daily momentum returns are modeled as

\begin{align}
r_d^{WML} &= \varepsilon_d, \\
\varepsilon_d &= \sigma_d z_d,
\qquad z_d \sim \mathcal{N}(0,1), \\
\sigma_d^2
&=
\omega
+
\alpha \varepsilon_{d-1}^2
+
\beta \sigma_{d-1}^2.
\label{eq:garch}
\end{align}

The model captures volatility clustering by
allowing conditional variance to depend on
recent squared return innovations and its
own lagged value.

At each month-end, we re-estimate the model
using a rolling window of 756 daily momentum
returns, approximately three trading years.
The parameters are estimated by maximum
likelihood under conditionally Gaussian errors,
with the conditional mean fixed at zero. For numerical
stability, daily returns are multiplied by 100 before
estimation, and forecast variances are converted back to
decimal-return units by dividing by $100^2$. Fits with a
nonzero optimizer convergence flag are treated as unavailable
forecasts rather than accepted estimates.

We then generate forecasts of daily conditional
variance for the next 21 trading days.
The monthly volatility forecast is constructed as

\begin{equation}
\widehat{\sigma}_{t+1|t}^{GARCH}
=
\sqrt{
\sum_{h=1}^{21}
\widehat{\sigma}_{t+h|t}^{2}
}.
\label{eq:garch_monthly}
\end{equation}

This approach incorporates the time-varying
conditional variance dynamics estimated by
GARCH while maintaining the same approximate
monthly forecasting horizon as the historical
estimator.

The historical and GARCH volatility forecasts are computed using
information available at the end of month $t$.
Their corresponding exposure multipliers are
applied to the momentum return realized during
month $t+1$, avoiding look-ahead bias.


\subsection{EWMA Volatility Forecast}

As a second extension, we consider an exponentially weighted
moving average (EWMA) estimator of momentum volatility.
Unlike the historical estimator, which assigns equal weight
to the preceding 126 daily observations, EWMA places
greater weight on recent squared returns while retaining
information from earlier observations.

Let $r_d^{WML}$ denote the daily momentum return.
The EWMA conditional variance estimate evolves according to

\begin{equation}
    v_d = \lambda v_{d-1}
    + (1-\lambda)(r_d^{WML})^2,
\end{equation}

where $\lambda \in (0,1)$ is the decay factor. Our baseline
specification sets $\lambda = 0.94$, consistent with the
conventional daily RiskMetrics parameterization.

The recursion is initialized using the mean squared return
over the first 126 available daily observations:

\begin{equation}
    v_{d_0} =
    \frac{1}{126}
    \sum_{i=1}^{126}(r_i^{WML})^2,
\end{equation}

where $d_0$ denotes the final observation in the
initialization window.

At each month-end $t$, we use the most recently updated
daily variance estimate to forecast volatility over the
following 21 trading days:

\begin{equation}
    \widehat{\sigma}^{EWMA}_{t+1|t}
    = \sqrt{21v_{d(t)}},
\end{equation}

where $d(t)$ denotes the final trading day of month $t$.

The forecast determines the exposure multiplier applied
to the subsequent month's momentum return, following
the same 12\% annualized volatility-targeting rule used
for the historical and GARCH strategies.

EWMA provides an intermediate specification between
the equally weighted historical estimator and the
parametrically estimated GARCH(1,1) model. It captures
volatility persistence through an exponentially decaying
weighting scheme without requiring repeated
maximum-likelihood estimation.

To assess sensitivity to the decay factor, we also
evaluate $\lambda \in \{0.90, 0.94, 0.97, 0.99\}$.
All specifications use the same underlying momentum
returns and are compared over identical monthly
observations within each evaluation period.

\subsection{Performance Evaluation}

We evaluate unmanaged momentum and all three
volatility-managed strategies using annualized
mean returns, annualized volatility, Sharpe
ratios, maximum drawdowns, and average exposure.

Annualized mean return and volatility are
calculated from monthly strategy returns as

\begin{align}
\mu_{\mathrm{ann}}
&= 12\,\overline{r}, \\
\sigma_{\mathrm{ann}}
&= \sqrt{12}\,s_r,
\end{align}

where $\overline{r}$ and $s_r$ are the sample
mean and standard deviation of monthly returns.

Sharpe ratios are computed as

\begin{equation}
SR
=
\frac{\mu_{\mathrm{ann}}}
{\sigma_{\mathrm{ann}}},
\end{equation}

using a zero risk-free rate for consistency
across strategies.

The baseline performance comparisons reported in Sections 4.1--4.6 are gross of transaction costs, financing expenses, and short-selling costs. These results characterize the theoretical performance of the volatility-management rules before implementation frictions. Sections 4.7--4.9 extend the analysis to stock-level transaction costs, leverage constraints, financing charges, and statistical inference on net strategy returns. Short-borrow fees and other unmodeled implementation costs remain excluded.


\subsection{Volatility Forecast Evaluation}

To evaluate the predictive accuracy of the three volatility
estimators, we compare their forecasts with the realized
volatility of the momentum portfolio during the subsequent
calendar month.

Realized monthly volatility is calculated from the daily
momentum returns observed within month $t+1$:

\begin{equation}
\sigma_{t+1}^{realized}
=
\sqrt{
\sum_{d \in \mathcal{M}_{t+1}}
(r_d^{WML})^2
},
\label{eq:realized_vol}
\end{equation}

where $\mathcal{M}_{t+1}$ denotes the set of trading days
in month $t+1$.

For each forecasting method
$j \in \{\mathrm{HIST},\mathrm{EWMA},\mathrm{GARCH}\}$,
we define the forecast error as

\begin{equation}
e_{t+1}^{j}
=
\widehat{\sigma}_{t+1|t}^{j}
-
\sigma_{t+1}^{realized}.
\end{equation}

We assess forecast accuracy using four measures:
mean absolute error (MAE), root mean squared error
(RMSE), mean forecast bias, and the correlation
between forecast and realized volatility.

These measures are defined as

\begin{align}
MAE_j
&=
\frac{1}{T}
\sum_{t=1}^{T}|e_t^j|,
\\[6pt]
RMSE_j
&=
\sqrt{
\frac{1}{T}
\sum_{t=1}^{T}(e_t^j)^2
},
\\[6pt]
Bias_j
&=
\frac{1}{T}
\sum_{t=1}^{T}e_t^j,
\\[6pt]
Corr_j
&=
Corr(
\widehat{\sigma}_{t}^{j},
\sigma_t^{realized}
).
\end{align}

Lower MAE and RMSE indicate more accurate volatility
forecasts, while a bias closer to zero indicates
less systematic overprediction or underprediction.
Higher correlation indicates a stronger association
between forecast and realized volatility.


\subsection{Statistical Significance of Forecast Improvements}

To determine whether differences in volatility forecast
accuracy are statistically significant, we compare the
squared forecast errors of the historical, EWMA, and
GARCH estimators.

For two competing forecasting methods $A$ and $B$,
we define the monthly loss differential as

\begin{equation}
    d_t^{A,B}
    =
    (e_t^A)^2 - (e_t^B)^2,
\end{equation}

where $e_t^A$ and $e_t^B$ denote the corresponding
volatility forecast errors. A positive loss differential
indicates that method $B$ produces a smaller squared
forecast error than method $A$.

We test the null hypothesis of equal expected squared
forecast losses,

\begin{equation}
    H_0: \mathbb{E}[d_t^{A,B}] = 0,
\end{equation}

against the two-sided alternative

\begin{equation}
    H_1: \mathbb{E}[d_t^{A,B}] \neq 0.
\end{equation}

The test statistic is

\begin{equation}
    t_{\mathrm{HAC}}
    =
    \frac{\overline{d}^{A,B}}
    {\widehat{\mathrm{SE}}_{\mathrm{HAC}}
    (\overline{d}^{A,B})},
\end{equation}

where the denominator is a heteroskedasticity- and
autocorrelation-consistent standard error of the
sample mean loss differential.

Following Newey and West (1987), we use HAC standard
errors with three monthly lags to account for potential
heteroskedasticity and serial dependence.

Our baseline analysis considers three pairwise comparisons:
historical volatility versus GARCH, historical volatility
versus EWMA, and GARCH versus EWMA.

We conduct these tests separately over the historical
sample (August 1929--December 2011) and the extension
sample (January 2012--December 2025).

We additionally assess forecast-loss differences across
alternative GARCH specifications and EWMA decay factors.
Because these robustness exercises involve multiple
comparisons, we distinguish individual test significance
from conclusions that remain significant after adjustment
for multiple testing.

Importantly, these tests evaluate statistical forecast
accuracy rather than differences in realized portfolio
performance. A more accurate volatility forecast does
not necessarily produce a higher Sharpe ratio when
used to determine portfolio exposure.

\subsection{Trading Frictions, Leverage Constraints, and Financing Costs}
\label{sec:implementation_method}
We complement the gross-return comparison with a stock-level, self-financing account simulation for the reconstructed winner-minus-loser (WML) portfolio. The simulator uses monthly holdings in the winner and loser deciles, applies the volatility-management exposure multiplier to their signed portfolio weights, and accounts for the trading needed to reach target positions. Trading costs are charged on absolute stock-level turnover, including rebalancing caused by changing constituents and drifted positions. Rather than deducting costs from an exogenously fixed portfolio value, the simulator solves for post-cost net asset value so that target positions and transaction costs are jointly self-financing. Monthly stock returns drive profit and loss, while ex-distribution price returns govern position drift; the difference is assigned to cash as a monthly approximation to distribution timing. The account does not simulate intramonth dividend reinvestment trades or delayed settlement of terminal delisting payoffs. Terminal delisting settlements and unresolved observations are handled explicitly by the underlying reconstruction.

Let $q_{i,t}$ denote the target signed weight of stock $i$ as a fraction of post-trading-cost NAV, $p_{i,t}$ its pretrade signed weight relative to pretrade NAV, and $c$ the proportional transaction-cost rate. The post-cost-to-pre-cost NAV ratio $x_t$ satisfies
\begin{equation}
 x_t+c\sum_i\left|x_tq_{i,t}-p_{i,t}\right|=1.
 \label{eq:self_financing_cost}
\end{equation}
The associated stock-level turnover is $\mathcal{T}_t=\sum_i|x_tq_{i,t}-p_{i,t}|$. Our baseline cost assumption is $c=10/10{,}000$, and the sensitivity analysis considers $0$, $5$, $10$, $20$, and $50$ basis points per unit traded. These assumptions represent stylized proportional trading fees rather than empirically estimated historical execution costs.

For leverage sensitivity, we replace the unconstrained exposure $L_t$ with $L_t^{(K)}=\min(L_t,K)$, where $K\in\{1,1.5,2\}$; the unconstrained rule is retained as a benchmark. We also consider annualized financing spreads $f\in\{0,100,300\}$ basis points on exposure exceeding one unit. The monthly financing deduction, measured relative to post-cost NAV, is
\begin{equation}
 F_t=\frac{f}{10{,}000\times12}\max\left(L_t^{(K)}-1,0\right).
 \label{eq:financing}
\end{equation}
The exposure multiplier scales the dollar-neutral WML portfolio rather than measuring gross absolute stock exposure. Because the underlying portfolio holds a long and a short decile, a multiplier of one corresponds to approximately two units of gross absolute stock exposure per unit of NAV. The financing convention charges only for the portion of the WML multiplier exceeding one; it does not model actual margin borrowing, short-sale proceeds, or security-borrowing expenses.

For gross portfolio return $R_t^{\mathrm{gross}}=\sum_iq_{i,t}r_{i,t}$, the modeled net account return is
\begin{equation}
 R_t^{\mathrm{net}}=x_t\left(1+R_t^{\mathrm{gross}}-F_t\right)-1.
 \label{eq:net_account_return}
\end{equation}
The financing charge is a simplified spread on leverage above one, not a comprehensive treatment of short-borrow fees, collateral remuneration, margin, market impact, or execution constraints. We evaluate the same 989 historical and 168 extension months as in the main reconstructed-portfolio comparison.

\subsection{Statistical Inference for Strategy Performance}
\label{sec:bootstrap_method}
To distinguish observed Sharpe-ratio differences from sampling uncertainty, we apply a paired circular moving-block bootstrap to the monthly net return series. In each replication, we sample starting months uniformly and concatenate circular blocks of consecutive observations until the original sample length is reached. The same resampled month indices are used for unmanaged momentum, historical volatility management, EWMA, and GARCH, preserving their contemporaneous return dependence while approximating short-run serial dependence.

For strategies $A$ and $B$, our statistic is $\Delta\widehat{SR}_{A,B}=\widehat{SR}_A-\widehat{SR}_B$, where the Sharpe ratios are computed from annualized arithmetic mean returns and sample standard deviations with a zero risk-free benchmark. We also examine annualized mean-return differences. The primary specification uses $5{,}000$ bootstrap replications, six-month blocks, 95\% percentile confidence intervals, and the 10-basis-point transaction-cost scenario with zero financing costs and no leverage cap. Sensitivity analyses use three- and twelve-month blocks and a one-times leverage cap. Reported approximate two-sided $p$-values are obtained by centering the bootstrap differences at their observed value and comparing the absolute centered differences with the absolute observed statistic. We report multiple-testing sensitivity using a Bonferroni threshold of $0.05/6$ within each family of six pairwise strategy comparisons. Percentile intervals and centered-bootstrap $p$-values are not guaranteed to give identical significance decisions in finite samples.

The procedure conditions on the estimated strategies and observed return histories: it does not re-estimate GARCH parameters or reconstruct momentum holdings inside each resample. Accordingly, inference concerns the realized strategy return series rather than the full uncertainty of the model-selection and portfolio-construction process.



\section{Empirical Results}
\label{sec:results}
\subsection{Replication of Published Results}
\label{sec:replication}

\begin{table}[htbp]
\centering
\caption{Replication of Barroso and Santa-Clara (2015): Different Sample Starts}
\label{tab:replication}
\small
\begin{tabular}{lrrrr}
\toprule
& \multicolumn{2}{c}{Unmanaged WML}
& \multicolumn{2}{c}{Managed WML} \\
\cmidrule(lr){2-3}\cmidrule(lr){4-5}
Statistic & Published & Replicated & Published & Replicated \\
\midrule
Annualized mean return (\%) & 14.46 & 14.80 & 16.50 & 16.67 \\
Annualized volatility (\%) & 27.53 & 27.45 & 16.95 & 16.78 \\
Sharpe ratio & 0.53 & 0.54 & 0.97 & 0.99 \\
Worst monthly return (\%) & -78.96 & -77.66 & -28.40 & -24.16 \\
Skewness & -2.47 & -2.38 & -0.42 & -0.35 \\
Excess kurtosis & 18.24 & 18.01 & 2.68 & 2.00 \\
Maximum drawdown (\%) & -96.69 & -95.56 & -45.20 & -43.83 \\
\bottomrule
\end{tabular}

\vspace{0.5em}
\begin{minipage}{0.98\textwidth}
\footnotesize
\textit{Notes:} Published statistics are from Barroso and
Santa-Clara (2015), Table 3 and the accompanying discussion,
for March 1927--December 2011. Replicated statistics use
Kenneth French's value-weighted momentum decile portfolios
for July 1927--December 2011. The start dates and early daily-return
sources differ, so this is not an identical-sample replication.
Managed strategies target
12\% annualized volatility using the preceding 126 daily
returns. Returns are gross of trading and financing costs.
\end{minipage}
\end{table}



Table~\ref{tab:replication} compares our replication of the
volatility-managed momentum strategy with the results reported
by Barroso and Santa-Clara (2015). The replication closely
reproduces the principal empirical findings of the original
study, although the reported statistics are not identical.

For unmanaged momentum, our replication produces an annualized
mean return of 14.80\%, annualized volatility of 27.45\%, and
a Sharpe ratio of 0.54. These estimates are close to the
published values of 14.46\%, 27.53\%, and 0.53, respectively.
The unmanaged portfolio also exhibits substantial downside
risk, with a worst monthly return of -77.66\%, negative
skewness of -2.38, and excess kurtosis of 18.01.

Applying the historical volatility-management rule improves
the risk-adjusted performance and return distribution of the
momentum strategy. In our replication, the annualized Sharpe
ratio increases from 0.54 to 0.99, compared with an increase
from 0.53 to 0.97 in the original study. Annualized volatility
declines from 27.45\% to 16.78\%, while the maximum drawdown
improves from -95.56\% to -43.83\%. The managed strategy also
exhibits substantially less negative skewness and lower excess
kurtosis.

The remaining numerical differences reflect limitations in
the comparability of the underlying datasets and sample
periods. The original study reports results from March 1927
through December 2011 and incorporates historical daily
momentum portfolios from Daniel and Moskowitz for the early
sample. 
Our replication instead uses Kenneth French's
value-weighted momentum decile portfolios throughout
and adopts July 1927 as the beginning of the
evaluation window. Although sufficient daily
observations are available to construct the first
volatility-managed return in May 1927, we retain
the July starting date for our baseline replication.
A sensitivity check beginning in May 1927 produces
virtually unchanged performance statistics.
 Differences
in data vintages and portfolio return histories may therefore
contribute to discrepancies, particularly in tail-sensitive
statistics.

Overall, the close agreement in average returns, volatility,
Sharpe ratios, and distributional characteristics supports
the validity of our implementation of the original
volatility-management methodology. Having established this
baseline, we turn to our principal extension: comparing the
historical volatility estimator with a rolling GARCH(1,1)
forecast using a consistently reconstructed CRSP momentum
portfolio.

\subsection{Portfolio Performance}


\begin{table}[htbp]
\centering
\caption{Momentum Portfolio Performance:
Historical and Extension Samples}
\label{tab:portfolio_performance}
\small
\begin{tabular}{lrrrrr}
\toprule
Strategy & Return (\%) & Vol. (\%) & Sharpe
& Max DD (\%) & Exposure \\
\midrule
\multicolumn{6}{l}{\textit{August 1929--December 2011}} \\
Unmanaged momentum
    & 14.10 & 27.65 & 0.510 & -95.37 & 1.000 \\
Historical volatility (126-day)
    & 16.68 & 17.32 & 0.963 & -42.02 & 0.912 \\
EWMA ($\lambda=0.94$)
    & 18.06 & 17.76 & 1.017 & -35.43 & 0.979 \\
GARCH(1,1)
    & 15.77 & 16.67 & 0.946 & -34.15 & 0.881 \\
\midrule
\multicolumn{6}{l}{\textit{January 2012--December 2025}} \\
Unmanaged momentum
    & 4.63 & 27.15 & 0.171 & -63.29 & 1.000 \\
Historical volatility (126-day)
    & 4.20 & 12.47 & 0.337 & -22.51 & 0.522 \\
EWMA ($\lambda=0.94$)
    & 4.56 & 13.30 & 0.343 & -24.68 & 0.562 \\
GARCH(1,1)
    & 4.13 & 12.99 & 0.318 & -25.86 & 0.530 \\
\bottomrule
\end{tabular}

\vspace{0.2cm}
\begin{minipage}{0.95\textwidth}
\footnotesize
\textit{Notes:} Gross returns, excluding transaction
and financing costs. Annualized returns are arithmetic
monthly means multiplied by 12. Sharpe ratios assume
a zero risk-free rate. Managed strategies target 12\%
annualized volatility without leverage caps.
All strategies use identical months within each sample.
\end{minipage}
\end{table}




Table~\ref{tab:portfolio_performance} compares the
performance of unmanaged momentum and three
volatility-managed strategies based on historical
volatility, EWMA, and GARCH(1,1) forecasts.

During the historical period (August 1929--December
2011), unmanaged momentum generates an annualized
mean return of 14.10\%, annualized volatility of
27.65\%, and a Sharpe ratio of 0.51. Its maximum
drawdown reaches 95.37\%, illustrating the substantial
downside risk of an unscaled momentum strategy.

All three volatility-management approaches substantially
improve risk-adjusted performance. Historical volatility
management increases the Sharpe ratio to 0.963,
while GARCH-based management achieves 0.946.
EWMA delivers the highest Sharpe ratio of 1.017,
with an annualized mean return of 18.06\% and
realized volatility of 17.76\%.

The EWMA strategy also reduces the maximum drawdown
to 35.43\%, compared with 42.02\% for the historical
estimator. GARCH produces a slightly smaller maximum
drawdown of 34.15\%, alongside the lowest annualized
volatility among the managed strategies.

The extension period (January 2012--December 2025)
reveals a substantial deterioration in unmanaged
momentum performance. Its annualized mean return
declines to 4.63\%, while volatility remains elevated
at 27.15\%, producing a Sharpe ratio of only 0.171.

Volatility management continues to improve
risk-adjusted performance in this more recent sample.
Historical volatility management achieves a Sharpe
ratio of 0.337, GARCH achieves 0.318, and EWMA
again records the highest value at 0.343.

However, the differences among managed strategies
are considerably smaller in the extension period.
EWMA produces annualized volatility of 13.30\%,
compared with 12.47\% for historical volatility
management and 12.99\% for GARCH. The historical
estimator also produces the smallest maximum drawdown
of 22.51\%, compared with 24.68\% for EWMA and
25.86\% for GARCH.

An important difference between the two periods
concerns average portfolio exposure. During the
historical sample, average exposure ranges from
0.881 for GARCH to 0.979 for EWMA. In the
extension period, exposure declines substantially
across all three managed strategies, ranging from
0.522 to 0.562.

These findings support the central conclusion of
Barroso and Santa-Clara (2015): volatility management
can substantially improve the risk-adjusted
performance of momentum strategies.

Among the three forecasting methods, EWMA produces the highest realized gross Sharpe ratio in both periods. Nevertheless, the comparisons presented in Table 2 are descriptive and do not, by themselves, establish statistically significant differences in risk-adjusted performance. Moreover, the strategies exhibit different realized volatility and exposure levels, and the results in this table exclude trading and financing costs. Consequently, the observed advantage of EWMA should not be interpreted as conclusive evidence of superior economic performance. Sections 4.7 and 4.8 subsequently examine the effects of transaction costs, leverage constraints, and financing expenses, while Section 4.9 formally evaluates differences in net Sharpe ratios using paired block-bootstrap inference.


\subsection{Volatility Forecast Accuracy}


\begin{table}[htbp]
\centering
\caption{Out-of-Sample Volatility Forecast Accuracy}
\label{tab:forecast_accuracy}
\small
\begin{tabular}{lrrrr}
\toprule
Estimator & MAE & RMSE & Bias & Correlation \\
\midrule
\multicolumn{5}{l}{\textit{August 1929--December 2011}} \\
Historical volatility (126-day)
    & 0.01730 & 0.02920 & 0.00077 & 0.723 \\
EWMA ($\lambda=0.94$)
    & 0.01453 & 0.02488 & -0.00139 & 0.803 \\
GARCH(1,1)
    & 0.01535 & 0.02507 & 0.00118 & 0.797 \\
\midrule
\multicolumn{5}{l}{\textit{January 2012--December 2025}} \\
Historical volatility (126-day)
    & 0.02269 & 0.03155 & 0.00234 & 0.662 \\
EWMA ($\lambda=0.94$)
    & 0.02064 & 0.02945 & -0.00090 & 0.713 \\
GARCH(1,1)
    & 0.02014 & 0.02909 & 0.00125 & 0.713 \\
\bottomrule
\end{tabular}

\vspace{0.2cm}
\begin{minipage}{0.95\textwidth}
\footnotesize
\textit{Notes:} Forecasts formed at month-end are
compared with the following month's realized
volatility, measured as the square root of the sum
of squared daily WML returns. All estimators
are evaluated on identical forecast origins within
each sample. Bias is defined as forecast minus
realized volatility.
\end{minipage}
\end{table}



Table~\ref{tab:forecast_accuracy} compares the
predictive accuracy of the 126-day historical
volatility estimator, EWMA, and GARCH(1,1).
Forecasts are evaluated against realized monthly
momentum volatility using mean absolute error
(MAE), root mean squared error (RMSE), forecast
bias, and correlation.

During the historical period, both EWMA and
GARCH improve forecast accuracy relative to the
historical estimator. The historical estimator
produces an RMSE of 0.02920, compared with
0.02488 for EWMA and 0.02507 for GARCH.
These correspond to reductions of approximately
14.8\% and 14.1\%, respectively.

EWMA also achieves the lowest MAE (0.01453)
and the highest correlation with realized
volatility (0.803), narrowly outperforming
GARCH on both measures. These results suggest
that assigning greater weight to recent squared
returns improves the ability to anticipate
changes in momentum volatility.

However, EWMA does not outperform the
historical estimator on every metric. Its
mean forecast bias is -0.00139, indicating
modest average underprediction, compared with
positive biases of 0.00077 and 0.00118 for
historical volatility and GARCH, respectively.
The historical estimator therefore exhibits
the smallest absolute bias in this period.

In the extension period, the forecasting
advantages of EWMA and GARCH remain present
but are smaller. GARCH achieves the lowest
RMSE (0.02909) and MAE (0.02014), followed
closely by EWMA (0.02945 and 0.02064).
Both improve upon the historical estimator,
which records an RMSE of 0.03155 and an
MAE of 0.02269.

EWMA exhibits the smallest absolute forecast
bias in the extension period (-0.00090),
while EWMA and GARCH achieve nearly
identical correlations with realized
volatility, approximately 0.713.

Overall, the results indicate that both
EWMA and GARCH improve statistical volatility
forecasting relative to the equally weighted
historical estimator. Their relative ranking,
however, depends on the evaluation period
and the chosen accuracy measure.

Notably, the more sophisticated GARCH model
does not consistently outperform the simpler
EWMA estimator. During the historical period,
EWMA achieves slightly lower forecast errors,
while GARCH has a marginal advantage in the
extension period.

These findings also highlight the distinction
between statistical forecasting performance
and economic outcomes. Although EWMA records
the highest realized Sharpe ratio in both
periods, its forecast accuracy is not
uniformly superior across samples.

Portfolio performance depends not only on
the magnitude of volatility forecast errors,
but also on their timing, the resulting
exposure adjustments, and subsequent
momentum returns.


\subsection{Statistical Significance}

\begin{table}[htbp]
\centering
\caption{HAC Tests of Volatility Forecast-Loss Differences}
\label{tab:forecast_significance}
\small
\begin{tabular}{lrrrr}
\toprule
Comparison ($A-B$) & Months & Mean Loss Diff.
& HAC $t$ & $p$-value \\
\midrule
\multicolumn{5}{l}{\textit{August 1929--December 2011}} \\
Historical -- GARCH
    & 989 & 0.000224 & 2.064 & 0.0391 \\
Historical -- EWMA
    & 989 & 0.000233 & 2.504 & 0.0123 \\
GARCH -- EWMA
    & 989 & 0.000009 & 0.279 & 0.7805 \\
\midrule
\multicolumn{5}{l}{\textit{January 2012--December 2025}} \\
Historical -- GARCH
    & 168 & 0.000149 & 1.008 & 0.3132 \\
Historical -- EWMA
    & 168 & 0.000128 & 0.890 & 0.3734 \\
GARCH -- EWMA
    & 168 & -0.000021 & -0.562 & 0.5738 \\
\bottomrule
\end{tabular}

\vspace{0.2cm}
\begin{minipage}{0.95\textwidth}
\footnotesize
\textit{Notes:} The loss differential equals the
squared forecast error of estimator $A$ minus
that of estimator $B$. Positive values favor $B$.
Tests use Newey--West HAC standard errors with
three monthly lags and two-sided $p$-values.
Reported $p$-values are unadjusted for multiple
comparisons.
\end{minipage}
\end{table}




Table~\ref{tab:forecast_significance} reports
HAC-adjusted tests comparing the squared volatility
forecast errors of the historical, EWMA, and
GARCH(1,1) estimators. Positive loss differentials
indicate that the second estimator in each
comparison achieves lower average squared
forecast errors.

During the historical period, both EWMA and
GARCH demonstrate statistically significant
improvements over the historical volatility
estimator when considered individually.

The historical-minus-GARCH loss differential
is 0.000224, with a HAC $t$-statistic of 2.064
and a two-sided $p$-value of 0.0391.
The historical-minus-EWMA differential is
slightly larger at 0.000233, with a
$t$-statistic of 2.504 and a $p$-value
of 0.0123.

These results provide evidence that both
conditional volatility approaches improve
squared-error forecast accuracy relative
to the equally weighted historical estimator
during the 1929--2011 sample.

However, the difference between EWMA and
GARCH is not statistically significant.
The GARCH-minus-EWMA loss differential is
only 0.000009, with a $t$-statistic of
0.279 and a $p$-value of 0.7805.
Consequently, we cannot establish that
either method provides more accurate
volatility forecasts than the other.

During the 2012--2025 extension, neither
EWMA nor GARCH produces a statistically
significant reduction in squared forecast
loss relative to the historical estimator.
The historical-minus-EWMA comparison
has a $p$-value of 0.3734, while the
historical-minus-GARCH comparison has
a $p$-value of 0.3132.

Similarly, the difference between GARCH
and EWMA remains statistically
insignificant ($p=0.5738$).

The shorter extension sample contains
168 monthly observations, compared with
989 in the historical period. This
difference in sample size may reduce
statistical power, making smaller
forecasting improvements more difficult
to detect. Nevertheless, differences
in statistical significance across
samples do not by themselves establish
a structural change in forecasting
performance.

Because the three baseline pairwise tests
within each sample are related, their
unadjusted $p$-values should be interpreted
with appropriate caution. Applying the
Holm correction to the three historical
comparisons retains significance for
historical versus EWMA, but not for
historical versus GARCH.

Overall, the results suggest that EWMA
provides a statistically supported
improvement over the historical estimator
in the historical sample, while GARCH
also exhibits an improvement under
unadjusted inference. Neither method
demonstrates a statistically established
forecasting advantage over the other.

Importantly, these tests concern
volatility forecast accuracy rather
than portfolio performance. Although
EWMA achieves the highest realized
Sharpe ratio in both periods, the
forecast-loss tests do not establish
that its investment performance is
statistically superior.



\subsection{Robustness Analysis: Innovation Distributions and Estimation Windows}
\label{subsec:garch_robustness}

To assess whether our GARCH results depend on the choice
of innovation distribution or estimation-window length,
we evaluate eight GARCH(1,1) specifications combining
Gaussian and Student-$t$ innovations with rolling windows
of 252, 504, 756, and 1,260 daily observations.

Each specification is estimated using the same reconstructed
daily WML returns and generates a 21-trading-day-ahead
volatility forecast at each month-end. These forecasts
determine next-month portfolio exposure according to
the volatility-targeting rule established in
Section 3.

The robustness analysis covers May 1931--December 2011
(968 months) and January 2012--December 2025
(168 months). The historical sample begins later than
the main comparison in Section 4.2 because the
1,260-day specification requires a longer initial
estimation period. All eight specifications are
therefore evaluated over identical monthly observations
within each period.

The Gaussian GARCH(1,1) model estimated using
756 daily observations serves as our benchmark.
Its forecasts exactly reproduce those of the
corresponding baseline specification, confirming
consistency between the main analysis and the
robustness implementation.


\subsubsection{Volatility Forecast Accuracy}

Table~\ref{tab:garch_robustness_accuracy} compares forecast
accuracy using mean absolute error (MAE), root mean squared
error (RMSE), and the correlation between forecast and realized
volatility. Realized monthly volatility is measured as the square
root of the sum of squared daily WML returns within the month.

\begin{table}[htbp]
\centering
\caption{GARCH Robustness: Volatility Forecast Accuracy}
\label{tab:garch_robustness_accuracy}
\small
\begin{tabular}{lrrrrrr}
\hline
& \multicolumn{3}{c}{1931--2011}
& \multicolumn{3}{c}{2012--2025} \\
\cline{2-4}\cline{5-7}
Specification & MAE & RMSE & Corr.
              & MAE & RMSE & Corr. \\
\hline
Gaussian 252  & 0.01576 & 0.02686 & 0.773
              & 0.02240 & 0.03172 & 0.680 \\
Gaussian 504  & 0.01528 & 0.02560 & 0.788
              & 0.02032 & 0.02907 & 0.714 \\
Gaussian 756  & 0.01518 & \textbf{0.02485} & 0.800
              & 0.02014 & 0.02909 & 0.713 \\
Gaussian 1260 & \textbf{0.01506} & 0.02496 & 0.797
              & 0.02062 & 0.02928 & 0.710 \\
Student-$t$ 252  & 0.01601 & 0.02865 & 0.754
                 & 0.02142 & 0.03046 & 0.697 \\
Student-$t$ 504  & 0.01566 & 0.02670 & 0.770
                 & 0.02027 & 0.02903 & 0.711 \\
Student-$t$ 756  & 0.01535 & 0.02528 & 0.791
                 & \textbf{0.02011} & 0.02896 & 0.712 \\
Student-$t$ 1260 & 0.01522 & 0.02511 & 0.794
                 & 0.02032 & \textbf{0.02874} & 0.717 \\
\hline
\end{tabular}
\begin{flushleft}
\footnotesize
\textit{Notes:} Forecasts are generated at month-end and compared
with realized volatility in the subsequent month. Lower MAE and
RMSE indicate greater forecast accuracy. Corr. denotes the
forecast--realized volatility correlation.
\end{flushleft}
\end{table}

In the historical sample, the Gaussian GARCH model estimated over
756 observations achieves the lowest RMSE (0.02485) and highest
forecast--realized volatility correlation (0.800). The 1,260-day
Gaussian specification achieves a slightly lower MAE, suggesting
that the preferred estimation window depends on the loss function.
More generally, longer estimation windows improve forecast
accuracy relative to the shortest windows, while Student-$t$
innovations do not consistently improve on Gaussian innovations.

The extension sample presents a somewhat different pattern.
Student-$t$ innovations become more competitive: the 756-day
Student-$t$ model achieves the lowest MAE (0.02011), while the
1,260-day Student-$t$ model achieves the lowest RMSE (0.02874).
However, the improvements relative to the Gaussian 756-day
benchmark are small.

To assess statistical significance, we calculate the monthly
squared-error loss differential

\begin{equation}
    d_t =
    \left(\widehat{\sigma}^{G756}_t-\sigma_t\right)^2
    -
    \left(\widehat{\sigma}^{j}_t-\sigma_t\right)^2,
\end{equation}

where $G756$ denotes the Gaussian 756-day benchmark and $j$
denotes an alternative specification. Positive values indicate
lower squared forecast error for the alternative. We test whether
the mean differential differs from zero using a constant-only
regression with heteroskedasticity- and autocorrelation-consistent
(HAC) standard errors and three lags.

\begin{table}[htbp]
\centering
\caption{HAC Tests of Forecast-Loss Differences}
\label{tab:garch_robustness_hac}
\small
\begin{tabular}{lrrrr}
\hline
& \multicolumn{2}{c}{1931--2011}
& \multicolumn{2}{c}{2012--2025} \\
\cline{2-3}\cline{4-5}
Alternative & $t$-stat. & $p$-value
            & $t$-stat. & $p$-value \\
\hline
Gaussian 252  & -2.084 & 0.037 & -1.892 & 0.059 \\
Gaussian 504  & -1.968 & 0.049 &  0.042 & 0.967 \\
Gaussian 1260 & -0.623 & 0.533 & -0.410 & 0.682 \\
Student-$t$ 252  & -1.830 & 0.067 & -1.482 & 0.138 \\
Student-$t$ 504  & -2.187 & 0.029 &  0.133 & 0.894 \\
Student-$t$ 756  & -1.086 & 0.277 &  0.393 & 0.694 \\
Student-$t$ 1260 & -0.871 & 0.384 &  0.719 & 0.472 \\
\hline
\end{tabular}
\begin{flushleft}
\footnotesize
\textit{Notes:} Tests compare each alternative against Gaussian
GARCH(1,1) estimated using 756 observations. Positive test
statistics favor the alternative. Two-sided $p$-values use
HAC standard errors with three lags. The seven comparisons
within each sample are not independent.
\end{flushleft}
\end{table}


In the historical sample, the Gaussian 756-day
benchmark achieves lower mean squared forecast error
than all seven alternatives. Three comparisons are
individually significant at the 5\% level:
Gaussian 252, Gaussian 504, and Student-$t$ 504.
However, none remains significant after applying
the Holm correction for multiple comparisons.
Thus, although the benchmark performs favorably,
the evidence is insufficient to establish its
superiority across the full set of specifications
after accounting for multiple testing.

In the extension sample, several Student-$t$
specifications achieve slightly lower forecast
errors than the Gaussian benchmark, but none
of the differences is statistically significant.
These findings suggest that the relative
forecasting performance of alternative GARCH
specifications is sensitive to the sample period
and evaluation criterion. They do not, however,
provide compelling evidence that changing the
innovation distribution or estimation window
consistently improves volatility forecasting.


\subsubsection{Portfolio Performance and Economic Implications}

Table~\ref{tab:garch_robustness_performance} reports the
annualized Sharpe ratio, annualized realized volatility, and
average portfolio exposure for each specification.

\begin{table}[htbp]
\centering
\caption{GARCH Robustness: Volatility-Managed Momentum Performance}
\label{tab:garch_robustness_performance}
\small
\begin{tabular}{lrrrrrr}
\hline
& \multicolumn{3}{c}{1931--2011}
& \multicolumn{3}{c}{2012--2025} \\
\cline{2-4}\cline{5-7}
Specification & Sharpe & Vol. & Exposure
              & Sharpe & Vol. & Exposure \\
\hline
Gaussian 252  & 0.978 & 17.33\% & 0.927
              & \textbf{0.377} & 12.90\% & 0.535 \\
Gaussian 504  & 0.943 & 16.93\% & 0.901
              & 0.302 & 13.08\% & 0.530 \\
Gaussian 756  & 0.938 & 16.64\% & 0.889
              & 0.318 & 12.99\% & 0.530 \\
Gaussian 1260 & 0.953 & 16.31\% & 0.879
              & 0.315 & 12.80\% & 0.525 \\
Student-$t$ 252  & \textbf{0.983} & 17.34\% & 0.929
                 & 0.372 & 12.95\% & 0.538 \\
Student-$t$ 504  & 0.934 & 16.94\% & 0.903
                 & 0.313 & 13.13\% & 0.535 \\
Student-$t$ 756  & 0.927 & 16.66\% & 0.889
                 & 0.318 & 13.08\% & 0.535 \\
Student-$t$ 1260 & 0.939 & 16.33\% & 0.876
                 & 0.314 & 12.88\% & 0.530 \\
\hline
\end{tabular}
\begin{flushleft}
\footnotesize
\textit{Notes:} Strategies target 12\% annualized volatility
using month-end forecasts and next-month WML returns.
Sharpe ratios are calculated using zero as the risk-free rate.
Exposure denotes the average portfolio weight. No leverage
cap or transaction costs are imposed.
\end{flushleft}
\end{table}

An important finding is that superior volatility forecasting does
not necessarily translate into superior risk-adjusted portfolio
performance. In the historical sample, the Gaussian 756-day
specification delivers the lowest forecast RMSE, but its Sharpe
ratio of 0.938 is below that of the Gaussian 252-day model (0.978)
and the Student-$t$ 252-day model (0.983). Similarly, in the
extension sample, Gaussian 252 days generates the highest Sharpe
ratio (0.377), despite having the highest forecast RMSE among
the eight specifications.

One possible explanation is that shorter estimation windows adjust more rapidly to changes in market conditions, thereby altering exposure during periods that are particularly important for realized portfolio returns. However, this mechanism has not been separately identified. Furthermore, the observed Sharpe-ratio differences among the eight alternative GARCH specifications have not been subjected to formal statistical tests. The paired block-bootstrap analysis in Section 4.9 evaluates the baseline historical volatility, EWMA, and GARCH strategies, but does not extend to comparisons among the alternative GARCH estimation windows and innovation distributions considered here. Accordingly, the specification-level performance rankings should be interpreted as descriptive rather than evidence of statistically significant differences.

Furthermore, the strategies do not achieve identical realized
volatility or average exposure. In the historical sample,
annualized realized volatility ranges from approximately 16.3\%
to 17.3\%, exceeding the nominal 12\% target. In the extension
sample, realized volatility is closer to the target, ranging
from approximately 12.8\% to 13.1\%. These differences should
be considered when interpreting the economic performance
comparisons.


Overall, the robustness analysis indicates that
the Gaussian 756-day specification provides a
reasonable forecasting benchmark, but its
performance is not uniformly superior across
all evaluation criteria. More importantly,
the results reinforce the central finding of
our GARCH extension: improvements in statistical
volatility forecast accuracy do not necessarily
translate into improvements in risk-adjusted
investment performance. The economic value of
a forecasting model depends not only on the
magnitude of its prediction errors, but also
on how those errors affect portfolio exposure
and subsequent realized returns.



\subsection{Robustness Analysis: EWMA Decay Factor}

To evaluate whether the performance of EWMA-based
volatility management depends on the choice of decay
factor, we consider four specifications with
$\lambda \in \{0.90, 0.94, 0.97, 0.99\}$.

The baseline value of $\lambda=0.94$ follows the
conventional daily RiskMetrics parameterization.
Lower decay factors assign greater weight to recent
squared returns and therefore allow estimated
volatility to respond more rapidly to new shocks.
Higher decay factors generate smoother, more
persistent volatility estimates.

All four specifications use the same reconstructed
daily momentum returns, 126-observation
initialization procedure, and 21-trading-day
forecast horizon. Comparisons are conducted over
identical monthly observations within each sample.

\begin{table}[htbp]
\centering
\caption{EWMA Robustness: Forecast Accuracy and
Portfolio Performance}
\label{tab:ewma_robustness}
\small
\begin{tabular}{lrrrr}
\toprule
Decay Factor & Sharpe & Vol. (\%) & RMSE
& Exposure \\
\midrule
\multicolumn{5}{l}{\textit{August 1929--December 2011}} \\
$\lambda=0.90$ & 1.010 & 18.44 & 0.02483 & 1.016 \\
$\lambda=0.94$ & 1.017 & 17.76 & 0.02488 & 0.979 \\
$\lambda=0.97$ & 1.004 & 17.22 & 0.02570 & 0.935 \\
$\lambda=0.99$ & 0.949 & 16.69 & 0.02853 & 0.869 \\
\midrule
\multicolumn{5}{l}{\textit{January 2012--December 2025}} \\
$\lambda=0.90$ & 0.336 & 14.06 & 0.02986 & 0.587 \\
$\lambda=0.94$ & 0.343 & 13.30 & 0.02945 & 0.562 \\
$\lambda=0.97$ & 0.342 & 12.65 & 0.02926 & 0.536 \\
$\lambda=0.99$ & 0.327 & 12.39 & 0.03159 & 0.504 \\
\bottomrule
\end{tabular}

\vspace{0.2cm}
\begin{minipage}{0.95\textwidth}
\footnotesize
\textit{Notes:} All strategies target 12\% annualized
volatility. Forecast RMSE is calculated against
the following month's realized volatility.
Volatility denotes annualized realized portfolio
volatility, and exposure is the average portfolio
multiplier. Returns exclude transaction and
financing costs.
\end{minipage}
\end{table}

Table~\ref{tab:ewma_robustness} reports the
forecast accuracy and portfolio performance
of the four EWMA specifications.

During the historical period, the baseline
$\lambda=0.94$ specification achieves the
highest Sharpe ratio of 1.017. The nearby
specifications with $\lambda=0.90$ and
$\lambda=0.97$ produce Sharpe ratios of
1.010 and 1.004, respectively, indicating
that the observed risk-adjusted performance
is relatively stable across moderate
changes in the decay factor.

The $\lambda=0.90$ specification achieves
the lowest historical forecast RMSE of
0.02483, marginally outperforming the
baseline value of 0.02488. However, the
difference is economically small and
does not translate into a higher
realized Sharpe ratio.

Increasing the decay factor to 0.99
produces a noticeably higher forecast
RMSE of 0.02853 and reduces the Sharpe
ratio to 0.949. This suggests that
placing excessive weight on older
observations may reduce the estimator's
responsiveness to changes in momentum
volatility.

In the extension sample, the baseline
specification again achieves the highest
Sharpe ratio of 0.343, although the
$\lambda=0.97$ specification performs
almost identically at 0.342.

The lowest extension-period forecast
RMSE is obtained with $\lambda=0.97$
(0.02926), compared with 0.02945
for the baseline specification.
As in the historical period, the
specification with $\lambda=0.99$
produces the lowest Sharpe ratio
among the four alternatives.

These findings indicate that the
baseline EWMA strategy does not
depend on a narrowly selected
decay parameter. Its performance
remains relatively stable for
$\lambda$ values between 0.90
and 0.97, while the slowest-decaying
specification performs less favorably.

Nevertheless, the relative rankings
of the specifications are descriptive
and do not establish statistically
significant differences in portfolio
performance.


We further compare squared volatility forecast
losses using HAC-adjusted tests with three
monthly lags, taking $\lambda=0.94$ as the
benchmark.

During the historical sample, neither
$\lambda=0.90$ ($p=0.9057$) nor
$\lambda=0.97$ ($p=0.1487$) differs
significantly from the baseline.
However, the $\lambda=0.99$ specification
produces significantly higher squared
forecast loss ($p=0.0079$).
This result remains significant after
applying the Holm correction across the
three alternative decay factors.

In the extension sample, none of the
alternative specifications differs
significantly from the baseline,
with $p$-values ranging from 0.3675
to 0.8158.

Taken together, these results support
the use of $\lambda=0.94$ as a reasonable
baseline rather than establishing it
as a universally optimal parameter.
The evidence favors moderate decay
factors over the slowest-decaying
specification in the historical sample,
but does not establish a statistically
superior decay factor in the
post-2011 period.

\subsection{Transaction-Cost-Adjusted Performance}
\label{sec:transaction_cost_results}
Table~\ref{tab:transaction_costs} presents the performance of the stock-level, self-financing account with a proportional trading cost of 10 basis points and no financing charge. Unlike the gross results in Table~2, these figures include costs associated with trading individual stock positions. The corresponding zero-cost simulations reproduce the baseline gross-return results within numerical precision. All four strategies remain evaluable over the 989 historical and 168 extension months, and the transaction-cost accounting identities pass the implementation checks.

\begin{table}[htbp]
\centering
\caption{Momentum performance after 10-basis-point stock-level transaction costs}
\label{tab:transaction_costs}
\small
\begin{tabular}{llrrrr}
\hline
Sample & Strategy & Mean (\%) & Vol. (\%) & Sharpe & Turnover \\
\hline
1929--2011 & Unmanaged & 12.241 & 27.619 & 0.443 & 1.541 \\
 & Historical & 15.000 & 17.254 & 0.869 & 1.375 \\
 & EWMA & 16.149 & 17.689 & 0.913 & 1.561 \\
 & GARCH & 14.101 & 16.618 & 0.849 & 1.369 \\
\hline
2012--2025 & Unmanaged & 2.838 & 27.117 & 0.105 & 1.490 \\
 & Historical & 3.267 & 12.458 & 0.262 & 0.775 \\
 & EWMA & 3.477 & 13.285 & 0.262 & 0.899 \\
 & GARCH & 3.127 & 12.977 & 0.241 & 0.831 \\
\hline
\end{tabular}
\begin{minipage}{0.94\linewidth}\footnotesize\emph{Notes:} Annualized arithmetic mean returns and volatility; Sharpe ratios use a zero risk-free benchmark. Turnover is average monthly sum of absolute stock-weight trades. Strategies are uncapped and assume zero financing spread. Figures are rounded from the validated leverage-financing experiment.\end{minipage}
\end{table}

Transaction costs reduce annualized mean returns, but they do not eliminate the historical risk-adjusted advantage of volatility management. Relative to the corresponding zero-cost simulations, the 10-basis-point assumption reduces historical annualized mean returns by approximately 1.86 percentage points for unmanaged momentum, 1.68 for historical volatility targeting, 1.91 for EWMA, and 1.67 for GARCH. In the extension sample, the corresponding reductions are approximately 1.79, 0.93, 1.08, and 1.00 percentage points. The lower modern trading-cost drag for the managed strategies is consistent with their lower average exposure and stock-level turnover. These comparisons do not establish historical execution feasibility: realized costs would vary with liquidity, market conditions, and short-selling requirements.

\subsection{Leverage Caps and Financing Costs}
\label{sec:leverage_results}
Table~\ref{tab:leverage_caps} compares uncapped volatility management with a one-times exposure cap, holding stock-level transaction costs at 10 basis points and the financing spread at zero. In the historical sample, the cap binds in approximately 40.4\%, 45.0\%, and 39.8\% of months for the historical, EWMA, and GARCH strategies, respectively. It reduces their Sharpe ratios from 0.869 to 0.779, 0.913 to 0.830, and 0.849 to 0.804. Thus, part of the historical performance improvement relies on the ability to scale positions above one unit of exposure. Nevertheless, even with the cap, all three strategies retain higher observed Sharpe ratios than unmanaged momentum (0.443).

\begin{table}[htbp]
\centering
\caption{Sensitivity of net Sharpe ratios to a one-times leverage cap}
\label{tab:leverage_caps}
\small
\begin{tabular}{llrrr}
\hline
Sample & Strategy & Uncapped & Cap at 1$\times$ & Cap binding (\%) \\
\hline
1929--2011 & Historical & 0.869 & 0.779 & 40.4 \\
 & EWMA & 0.913 & 0.830 & 45.0 \\
 & GARCH & 0.849 & 0.804 & 39.8 \\
\hline
2012--2025 & Historical & 0.262 & 0.262 & 0.6 \\
 & EWMA & 0.262 & 0.261 & 6.5 \\
 & GARCH & 0.241 & 0.240 & 1.8 \\
\hline
\end{tabular}
\begin{minipage}{0.94\linewidth}\footnotesize\emph{Notes:} All returns incorporate 10-basis-point stock trading costs and assume zero financing spread. Binding fractions measure the proportion of months in which unconstrained exposure exceeds one.\end{minipage}
\end{table}

The extension sample displays a different pattern. Uncapped average exposure is only about 0.522 for historical volatility targeting, 0.562 for EWMA, and 0.530 for GARCH. Consequently, the one-times cap binds in fewer than 7\% of months for any managed strategy, and the resulting Sharpe changes are negligible. Raising the cap to 1.5 or 2 times likewise has little effect in the modern sample. The contrast suggests that leverage flexibility contributes more to historical managed-momentum performance than to the more recent results.

Financing costs further reduce historical performance, but their effect is smaller than that of a strict one-times cap. At 300 annualized basis points on exposure above one, the historical uncapped Sharpe ratios fall from 0.869 to 0.850 (historical volatility), 0.913 to 0.885 (EWMA), and 0.849 to 0.833 (GARCH). The modern Sharpe ratios change very little under the same assumption because leverage above one is uncommon. These findings are conditional on the simplified financing convention in Equation~\eqref{eq:financing}; they do not include stock-borrow fees or other costs of maintaining a short leg.

\subsection{Bootstrap Inference and Performance Robustness}
\label{sec:bootstrap_results}
Table~\ref{tab:bootstrap_sharpe} reports the primary paired six-month circular block-bootstrap results for net returns with 10-basis-point transaction costs, no leverage cap, and zero financing spread. In the historical sample, all three volatility-managed strategies have significantly higher Sharpe ratios than unmanaged momentum: estimated differences range from 0.405 to 0.470, with all 95\% percentile confidence intervals strictly above zero. In the 2012--2025 extension, the corresponding estimated differences remain positive, approximately 0.136--0.158, but every confidence interval includes zero. Thus, the shorter modern sample provides insufficient evidence at the conventional 5\% level to establish a nonzero Sharpe improvement; this does not establish that the true benefit has disappeared.

\begin{table}[htbp]
\centering
\caption{Paired block-bootstrap differences in net annualized Sharpe ratios}
\label{tab:bootstrap_sharpe}
\small
\begin{tabular}{llrrr}
\hline
Sample & Comparison ($A-B$) & $\Delta SR$ & 95\% percentile CI & Approx. $p$ \\
\hline
1929--2011 & Historical -- Unmanaged & 0.426 & [0.246, 0.593] & 0.0002 \\
 & EWMA -- Unmanaged & 0.470 & [0.277, 0.648] & 0.0002 \\
 & GARCH -- Unmanaged & 0.405 & [0.221, 0.572] & 0.0002 \\
 & EWMA -- Historical & 0.044 & [-0.012, 0.097] & 0.1174 \\
 & GARCH -- Historical & -0.021 & [-0.078, 0.045] & 0.5003 \\
 & EWMA -- GARCH & 0.064 & [0.019, 0.110] & 0.0052 \\
\hline
2012--2025 & Historical -- Unmanaged & 0.158 & [-0.028, 0.340] & 0.1004 \\
 & EWMA -- Unmanaged & 0.157 & [-0.040, 0.317] & 0.0896 \\
 & GARCH -- Unmanaged & 0.136 & [-0.053, 0.279] & 0.1066 \\
 & EWMA -- Historical & -0.001 & [-0.108, 0.109] & 0.9930 \\
 & GARCH -- Historical & -0.021 & [-0.138, 0.089] & 0.7081 \\
 & EWMA -- GARCH & 0.021 & [-0.021, 0.065] & 0.3435 \\
\hline
\end{tabular}
\begin{minipage}{0.94\linewidth}\footnotesize\emph{Notes:} 5,000 paired circular block-bootstrap resamples with six-month blocks. The same monthly indices are applied to all strategies. Approximate two-sided $p$-values are based on centered bootstrap differences. Historical and extension samples contain 989 and 168 monthly observations, respectively.\end{minipage}
\end{table}

The historical EWMA--GARCH Sharpe difference is positive and its primary confidence interval excludes zero; the approximate $p$-value of 0.0052 is also below the Bonferroni threshold $0.05/6\approx0.0083$. However, this apparent advantage is not robust to imposing a one-times leverage cap. Neither EWMA nor GARCH exhibits a consistently significant Sharpe improvement over the historical volatility estimator. These results reinforce the distinction between improving volatility forecast accuracy and improving realized risk-adjusted portfolio performance.

Sensitivity tests use three-, six-, and twelve-month blocks and repeat the comparisons under a one-times leverage cap. In the historical sample, each managed strategy outperforms unmanaged momentum at the 5\% level in all six configurations, and all 18 corresponding comparisons survive the within-configuration Bonferroni threshold. The historical EWMA--historical comparison has percentile intervals excluding zero in three of six configurations, but none of its approximate $p$-values survives Bonferroni adjustment. EWMA--GARCH has intervals excluding zero in three of six configurations and Bonferroni-adjusted significance in two; all three significant percentile intervals occur without the leverage cap. GARCH--historical is not significant in any configuration. In the extension sample, none of the six pairwise comparisons has a 95\% percentile interval excluding zero in any of the six configurations. These are sensitivity checks on dependent estimates, not independent experimental replications. Moreover, confidence intervals and centered-bootstrap $p$-values can yield different decisions near the 5\% threshold.
\section{Robustness and Limitations}
\label{sec:limitations}


Our analysis provides evidence that volatility management
improves the risk-adjusted performance of momentum
strategies and that both EWMA and GARCH(1,1) can
improve volatility forecast accuracy relative to the
historical estimator. Nevertheless, several methodological
considerations affect the interpretation and
generalizability of these findings.

\subsection{Momentum Portfolio Reconstruction}

Unlike studies relying exclusively on a published
momentum factor, we reconstruct the WML portfolio
using CRSP stock-level data. This approach allows
us to evaluate volatility forecasts on a common
underlying return series, but introduces additional
data-processing assumptions.

In particular, historical security records can
contain missing returns, incomplete delisting
payoffs, and gaps in price observations.
Our primary daily reconstruction retains available
numeric return components while assigning zero
increments to certain unresolved observations.
These assumptions are explicitly identified
within the reconstruction pipeline and should
not be interpreted as observed market returns.

Although our reconstructed momentum series has
been compared with the Kenneth French momentum
factor as an external validation benchmark,
close agreement between the series does not
eliminate uncertainty surrounding individual
missing observations or terminal payoffs.

Additionally, our reconstruction uses a
separately identified CRSP CIZ eligibility
screen for recent observations where legacy
share and exchange classifications are
unavailable. Differences between the CIZ
screen and historical classifications may
affect the composition of the momentum
portfolio in the final part of the sample.


\subsection{Volatility Forecast Specification}

Our analysis compares three volatility forecasting
approaches with different assumptions and levels
of complexity. The historical estimator assigns
equal weight to the preceding 126 daily squared
returns, EWMA applies exponentially declining
weights with a baseline decay factor of
$\lambda=0.94$, and GARCH(1,1) models conditional
variance dynamics using parameters estimated
over a rolling window of 756 daily observations.

These specifications provide a consistent
framework for evaluating alternative volatility
forecasts, but their baseline parameters were
not optimized separately for each evaluation
period. The EWMA decay factor follows the
conventional daily RiskMetrics parameterization,
while the historical and GARCH estimation
windows reflect our chosen benchmark designs.

All three methods forecast volatility over an
approximate 21-trading-day horizon, whereas
realized volatility is calculated using the
actual number of trading days in each calendar
month. Consequently, part of the measured
forecast error may reflect differences between
the fixed forecasting horizon and the realized
evaluation period rather than deficiencies
in the underlying volatility model.

The EWMA specification assumes a constant
decay factor and does not explicitly model
a long-run variance level. Its responsiveness
to new volatility shocks depends on the
chosen value of $\lambda$. Our robustness
analysis evaluates decay factors of
0.90, 0.94, 0.97, and 0.99, finding
relatively stable portfolio Sharpe ratios
for the first three specifications but
weaker historical forecast accuracy
when $\lambda=0.99$.

Our baseline GARCH specification assumes
conditionally Gaussian innovations and a
zero conditional mean. The GARCH robustness
analysis relaxes the Gaussian assumption
by considering Student-$t$ innovations
and alternative rolling estimation windows
of 252, 504, 756, and 1,260 daily observations.

However, we do not investigate asymmetric
volatility models such as GJR-GARCH or
EGARCH, which could potentially capture
different responses to positive and negative
momentum shocks. We also retain a zero
conditional mean across all GARCH
specifications.

Finally, the three estimators differ in
their computational requirements. EWMA
updates recursively without repeated
parameter estimation, while our GARCH
implementation requires rolling
maximum-likelihood estimation. Although
this distinction may matter operationally,
we do not formally quantify computational
costs or implementation efficiency.

These modeling choices and restrictions
remain potential directions for
future research.


\subsection{Trading Frictions and Leverage}

Our net-return experiments account for proportional stock-level transaction costs and stylized financing charges on exposure above one, and impose alternative exposure caps. The historical advantage of volatility management persists after 10-basis-point stock trading costs and under a one-times leverage cap, although performance declines relative to the gross unconstrained benchmark. The extension sample is less sensitive to leverage and financing because managed exposures average only about 0.52--0.56. These exercises do not capture time-varying spreads, market impact, borrow availability or fees, margin and collateral conventions, tax effects, or all historical execution constraints. The net results are therefore implementation sensitivities rather than fully investable backtests.



\subsection{Statistical Inference and Sample Dependence}

Several limitations should be considered when
interpreting the statistical significance of
our results.

First, the historical and extension samples
differ considerably in length. The main
comparison contains 989 historical monthly
observations and 168 extension observations,
while the GARCH robustness analysis uses
968 historical observations to accommodate
the longest estimation window.

The shorter extension period provides
less information for detecting differences
in forecast accuracy, potentially reducing
the statistical power of our tests.
Consequently, the absence of statistically
significant forecasting improvements
after 2011 should not be interpreted
as evidence that competing models have
identical predictive performance.

Second, our historical analysis is not
a genuinely independent post-publication
test of the original momentum-management
hypothesis. Although the 2012--2025
extension provides a more recent
evaluation, the two periods differ
in market conditions, momentum returns,
and average portfolio exposure.

Differences in statistical significance
across samples do not, by themselves,
establish a structural change in
forecasting performance.

Third, our HAC-adjusted tests evaluate
differences in mean squared volatility
forecast errors rather than differences
in realized portfolio performance.
The choice of squared-error loss
provides a consistent statistical
criterion, but alternative loss
functions could yield different
rankings of forecasting models.

Fourth, our analysis involves multiple
comparisons. The baseline analysis
compares three forecasting methods,
while the GARCH and EWMA robustness
exercises evaluate seven and three
alternative specifications, respectively.

We therefore distinguish unadjusted
pairwise significance from conclusions
that remain significant after applying
the Holm correction within each
comparison family. These adjustments
reduce the risk of false discoveries
but may also reduce statistical power.

Formal paired circular moving-block bootstrap comparisons now assess differences in realized net Sharpe ratios. These tests condition on the observed strategy return series and do not re-estimate forecasting models or reconstruct stock portfolios within bootstrap samples. Results can depend on block length and exposure constraints; percentile confidence intervals and centered-bootstrap p-values may also differ at conventional significance thresholds. We examine multiple comparisons using a within-family Bonferroni correction. The shorter extension sample limits statistical precision, and a lack of significance is not evidence of equal underlying performance. The historical sample is not a fully independent test of the original hypothesis.



\section{Conclusion}

This paper replicates and extends the
volatility-managed momentum framework
of Barroso and Santa-Clara (2015),
with particular emphasis on the
economic and statistical consequences
of alternative volatility forecasting
methods.

Our replication confirms the
importance of volatility management
for momentum strategies. During the
historical sample, scaling momentum
exposure using estimated volatility
substantially increases realized
Sharpe ratios and reduces maximum
drawdowns relative to the unmanaged
portfolio. These benefits persist
in the 2012--2025 extension,
although the underlying momentum
premium and the returns of all
strategies are considerably weaker.

Our methodological extension
compares a 126-day historical
volatility estimator with EWMA
and GARCH(1,1). Both conditional
volatility approaches generally
improve forecast accuracy relative
to the historical benchmark.
EWMA produces the lowest forecast
RMSE during the historical period,
whereas GARCH achieves a marginally
lower RMSE in the extension period.

Statistical tests provide evidence
that EWMA improves squared-error
forecast accuracy relative to the
historical estimator during the
historical sample, including after
adjusting for the three baseline
pairwise comparisons. GARCH also
demonstrates an improvement under
unadjusted inference. However,
neither estimator establishes a
statistically significant forecasting
advantage over the other, and
forecasting improvements are not
statistically significant in the
extension sample.

From an investment perspective,
EWMA achieves the highest observed
Sharpe ratio among the evaluated
strategies in both periods. Its
performance is also relatively
stable across decay factors
between 0.90 and 0.97.

These findings suggest that the
benefits of volatility-managed
momentum are robust to several
reasonable choices of volatility
estimator. They also highlight
an important distinction between
statistical forecasting accuracy
and realized investment
performance: more complex models
do not necessarily generate
economically superior portfolios.

Our implementation extensions show that the historical Sharpe advantage survives 10-basis-point stock-level transaction costs and a one-times leverage cap, although both assumptions reduce realized performance. Financing charges on leverage above one have a smaller modeled effect, particularly after 2011 when average managed exposure is well below one. Paired block-bootstrap tests find a statistically supported historical net Sharpe improvement for all three managed strategies across block lengths and leverage assumptions. In contrast, none of the managed-versus-unmanaged improvements is statistically significant in the 2012--2025 sample. Nor is there a robust, implementation-independent Sharpe advantage of EWMA or GARCH over the simple historical estimator. These results suggest that the principal contribution of volatility management lies in scaling momentum risk rather than in selecting a uniquely superior forecasting model.

Overall, our analysis reinforces
the relevance of volatility
management for momentum investing
while demonstrating that a
computationally simple EWMA
approach can remain competitive
with more sophisticated
GARCH-based forecasts.

\bibliographystyle{apalike}
\bibliography{references}

\end{document}